% !TeX root = ../../Book4.tex
% 纲要标题：### F.2 双范畴与张量积
% 纲要标签：b4:sec:F02
% 纲要细目（写作范围）：
% #### F.2.1 结合与单位约束
% #### F.2.2 环、双模与幺半范畴

\section{双范畴与张量积}
\label{b4:sec:F02}

双模可以把两个环联系起来，并用张量积连续复合。不过，\((P\otimes_C N)\otimes_B M\) 与 \(P\otimes_C(N\otimes_B M)\) 在通常的构造中是两个对象，联系它们的是指定的自然同构。要把这种复合纳入二维的范畴结构，就应把结合与单位同构也列为数据，并规定改括号的不同过程相容。这与\secref{b4:sec:monoidal-structure}中的幺半结构有相同的要求。

\subsection{结合与单位约束}
\label{b4:subsec:F02:01}

\begin{definition}\label{b4:def:F-bicategory}
给定一个对象类，对每对对象 \(X,Y\) 给定一个局部小范畴 \(\operatorname{Hom}_{\mathcal B}(X,Y)\)。其对象记作 \(f:X\rightarrow Y\)，称为 1态射；其态射记作 \(\alpha:f\Rightarrow f'\)，称为 2态射。以 \(\beta\cdot\alpha\) 表示这些范畴中的复合，以 \(1_f\) 表示恒等 2态射。再给定以下数据：
\begin{enumerate}
\item 对任意 \(X,Y,Z\)，有复合函子
\[
\operatorname{Hom}_{\mathcal B}(Y,Z)\times
\operatorname{Hom}_{\mathcal B}(X,Y)
\longrightarrow\operatorname{Hom}_{\mathcal B}(X,Z),
\]
它在对象与态射上的值分别记为 \(g*f\) 与 \(\beta*\alpha\)，次序均为先右后左。
\item 对每个 \(X\)，给定单位 1态射 \(1_X\in\operatorname{Hom}_{\mathcal B}(X,X)\)。
\item 对依次可复合的 \(f:X\rightarrow Y\)、\(g:Y\rightarrow Z\)、\(h:Z\rightarrow W\)，给定关于各变量自然的可逆 2态射
\[
\begin{aligned}
a_{h,g,f}&:(h*g)*f\Longrightarrow h*(g*f),\\
\ell_f&:1_Y*f\Longrightarrow f,
\qquad r_f:f*1_X\Longrightarrow f.
\end{aligned}
\]
\end{enumerate}
若对任意依次可复合的四个 1态射 \(f,g,h,k\)，满足五边形等式
\begin{equation}\label{b4:eq:F-bicategory-pentagon}
\begin{aligned}
a_{k,h,g*f}\cdot a_{k*h,g,f}
={}&(1_k*a_{h,g,f})\cdot a_{k,h*g,f}
\cdot(a_{k,h,g}*1_f),
\end{aligned}
\end{equation}
且对 \(f:X\rightarrow Y\)、\(g:Y\rightarrow Z\)，满足三角等式
\begin{equation}\label{b4:eq:F-bicategory-triangle}
(1_g*\ell_f)\cdot a_{g,1_Y,f}=r_g*1_f,
\end{equation}
则称这些数据为一个\term[shuang-fanchou]{双范畴}[bicategory]。\(a\) 称为结合约束，\(\ell,r\) 称为单位约束。
\end{definition}

五边形等式的左边对应两次改括号：
\[
\begin{aligned}
((k*h)*g)*f
&\xrightarrow{a_{k*h,g,f}}(k*h)*(g*f)\\
&\xrightarrow{a_{k,h,g*f}}k*(h*(g*f)).
\end{aligned}
\]
右边则经过另外两个中间对象：
\[
\begin{aligned}
((k*h)*g)*f
&\xrightarrow{a_{k,h,g}*1_f}(k*(h*g))*f\\
&\xrightarrow{a_{k,h*g,f}}k*((h*g)*f)\\
&\xrightarrow{1_k*a_{h,g,f}}k*(h*(g*f)).
\end{aligned}
\]
两条路径的起点与终点相同，五边形等式要求沿它们复合指定的结合约束，得到同一个 2态射。三角等式比较删除中间单位的两种方法：
\[
\begin{tikzcd}[column sep=large]
(g*1_Y)*f\ar[rr,"a_{g,1_Y,f}"]\ar[dr,"r_g*1_f"']
&&g*(1_Y*f)\ar[dl,"1_g*\ell_f"]\\
&g*f&
\end{tikzcd}
\]
图中箭头均为 \(\operatorname{Hom}_{\mathcal B}(X,Z)\) 中的态射，也就是双范畴的 2态射。约束的自然性还要求改括号与 2态射相容。例如，对 \(\alpha:f\Rightarrow f'\)、\(\beta:g\Rightarrow g'\)、\(\gamma:h\Rightarrow h'\)，有
\[
a_{h',g',f'}\cdot((\gamma*\beta)*\alpha)
=(\gamma*(\beta*\alpha))\cdot a_{h,g,f}.
\]
复合是函子，因而 2态射的交换律仍成立；各 Hom 范畴中的垂直复合也严格结合。若结合与单位约束全为恒等 2态射，就得到\cref{b4:def:F-strict-two-category}的严格 2范畴。用范畴、配对函子与相容同构组织这种结构的动机，可对照\cite[Introduction, p. 1]{Benabou1967Bicategories}。

\subsection{环、双模与幺半范畴}
\label{b4:subsec:F02:02}

第二卷第6.1--6.2节构造了平衡张量积及其剩余作用。先看含幺环同态 \(\varphi:A\rightarrow B\)。令 \(B\) 保留自身的左乘作用，右 \(A\) 作用取为 \(b\cdot a=b\varphi(a)\)，便得到双模 \({}_BB_A\)。它把左 \(A\) 模 \(X\) 送到左 \(B\) 模 \(B\otimes_A X\)，正是沿 \(\varphi\) 的扩张标量。因此规定从 \(A\) 到 \(B\) 的 1态射为左 \(B\)、右 \(A\) 双模：一般的 \({}_BM_A\) 也给出同方向的函子 \(M\otimes_A-\)。复合环同态所对应的双模怎样由张量积得到，将在\cref{b4:prop:F-ring-map-composition}中证明。

\begin{proposition}\label{b4:prop:F-bimodule-bicategory}
以含幺环为对象，以幺性 \((B,A)\) 双模为 \(\operatorname{Hom}_{\mathbf{Bimod}}(A,B)\) 的对象，以保持左右作用的群同态为其态射。对 \({}_BM_A:A\rightarrow B\)、\({}_CN_B:B\rightarrow C\)，令
\[
N*M=N\otimes_B M,
\qquad c(n\otimes m)a=(cn)\otimes(ma),
\]
并对双模同态 \(\phi:M\rightarrow M'\)、\(\psi:N\rightarrow N'\)，令
\[
(\psi*\phi)(n\otimes m)=\psi(n)\otimes\phi(m).
\]
取单位 1态射为 \({}_AA_A\)。对 \({}_DP_C:C\rightarrow D\)，结合与单位约束由
\[
\begin{aligned}
a_{P,N,M}&:(P\otimes_C N)\otimes_B M
\longrightarrow P\otimes_C(N\otimes_B M),\\
&(p\otimes n)\otimes m\longmapsto p\otimes(n\otimes m),\\
\ell_M&:B\otimes_B M\longrightarrow M,
\qquad b\otimes m\longmapsto bm,\\
r_M&:M\otimes_A A\longrightarrow M,
\qquad m\otimes a\longmapsto ma
\end{aligned}
\]
给出。这些数据组成双范畴 \(\mathbf{Bimod}\)。
\end{proposition}
\symidx{\(\mathbf{Bimod}\)：环、双模与双模同态构成的双范畴}
\begin{proof}
双模作用彼此相容，所以外侧的 \(C,A\) 作用保留平衡关系 \(nb\otimes m=n\otimes bm\)，得到所列 \((C,A)\) 双模结构。对同态，平衡关系的两边分别送到
\[
\psi(nb)\otimes\phi(m)
=\psi(n)b\otimes\phi(m)
=\psi(n)\otimes b\phi(m)
=\psi(n)\otimes\phi(bm).
\]
因此 \(\psi*\phi\) 良定义，且保持外侧作用。在纯张量上，先后施行两对同态等于张量各自的复合，恒等同态的张量也为恒等；纯张量生成张量积，故这给出复合函子。

结合公式对 \(C\) 与 \(B\) 的平衡关系分别由
\[
\begin{aligned}
pc\otimes(n\otimes m)&=p\otimes(cn\otimes m),\\
p\otimes(nb\otimes m)&=p\otimes(n\otimes bm)
\end{aligned}
\]
保持。反向公式 \(p\otimes(n\otimes m)\mapsto(p\otimes n)\otimes m\) 同样保持这两类关系，故由张量积泛性质得到互逆的群同态。它们保留左 \(D\)、右 \(A\) 作用，所以为双模同构。对 \(M,N,P\) 分别施行同态 \(\phi,\psi,\chi\) 时，两种次序都把 \((p\otimes n)\otimes m\) 送到 \(\chi(p)\otimes(\psi(n)\otimes\phi(m))\)，给出自然性。

单位公式由模作用给出，平衡关系使其良定义。逆公式分别为 \(m\mapsto1_B\otimes m\)、\(m\mapsto m\otimes1_A\)；例如 \(1_B\otimes bm=b\otimes m\)、\(ma\otimes1_A=m\otimes a\)，所以两次复合为恒等，并保留双模作用。同态保持作用，故这两项约束也自然。

再取 \({}_EQ_D:D\rightarrow E\)。五边形的两条路径都把
\[
((q\otimes_D p)\otimes_C n)\otimes_B m
\longmapsto q\otimes_D(p\otimes_C(n\otimes_B m)).
\]
三角形的两条路径在 \((n\otimes b)\otimes m\) 上分别给出 \(nb\otimes m\) 与 \(n\otimes bm\)，由 \(B\) 平衡关系相等。多重纯张量生成相应双模，因而这两项相容条件都是同态等式。
\end{proof}

在双范畴中，若存在 1态射 \(f:X\rightarrow Y\)、\(g:Y\rightarrow X\)，以及可逆 2态射 \(g*f\Rightarrow1_X\)、\(f*g\Rightarrow1_Y\)，便称对象 \(X,Y\) 等价。对于环与双模，这要求两个复合张量积分别同构于相应的单位双模。\cref{b4:prop:F-matrix-bimodule-equivalence}将证明矩阵环 \(M_n(k)\) 与 \(k\) 按此意义等价，并得到它们的左模范畴等价；当 \(n>1\) 时，这两个环本身并不同构。

\ProseParagraphBreak

在通常的张量积构造中，结合约束是上述双模同构，而非恒等映射；单位约束也需要删去张量中的环因子。因此这组具体数据给出双范畴。固定一个环 \(A\) 后，\(A\)-\(A\) 双模在 \(\otimes_A\) 下的幺半结构正是其中单对象的情形。

\begin{proposition}\label{b4:prop:F-one-object-bicategory}
设双范畴 \(\mathcal B\) 只有一个对象 \(O\)。范畴 \(\operatorname{Hom}_{\mathcal B}(O,O)\) 取 \(X\otimes Y=X*Y\)、单位对象 \(1_O\)，以及结合、单位约束 \(a,\ell,r\)，便成为幺半范畴。反过来，每个局部小幺半范畴 \((\mathcal C,\otimes,I,\alpha,\lambda,\rho)\) 都确定一个单对象双范畴：唯一的 Hom 范畴为 \(\mathcal C\)，复合函子为 \(\otimes\)，单位 1态射为 \(I\)，约束为 \(\alpha,\lambda,\rho\)。这两种构造逐项恢复原有结构数据，只须重新命名双范畴的唯一对象。
\end{proposition}
\begin{proof}
单对象时，所有 1态射均可复合，因而复合函子就是同一范畴上的双函子。五边形与三角等式分别变成\eqref{b4:eq:monoidal-pentagon}与\eqref{b4:eq:monoidal-triangle}，约束的自然性也不变。反向把幺半范畴的对象看成 1态射、态射看成 2态射，便得到完全相同的数据与公理。
\end{proof}

\begin{proposition}\label{b4:prop:F-ring-map-composition}
设 \(A,B,C\) 为含幺环，\(\varphi:A\rightarrow B\)、\(\psi:B\rightarrow C\) 为保持单位的环同态。令 \(B,C\) 保留各自的左乘作用，分别取右作用 \(b\cdot a=b\varphi(a)\)、\(c\cdot b=c\psi(b)\)，得到 \((B,A)\) 双模 \(B\) 与 \((C,B)\) 双模 \(C\)，则
\[
\vartheta:C\otimes_B B\longrightarrow C,
\qquad c\otimes b\longmapsto c\psi(b)
\]
为 \((C,A)\) 双模同构，其中右端的右 \(A\) 作用为 \(c\cdot a=c\psi(\varphi(a))\)。对每个幺性左 \(A\) 模 \(X\)，还有关于 \(X\) 自然的左 \(C\) 模同构
\[
\eta_X:C\otimes_B(B\otimes_A X)\longrightarrow C\otimes_A X,
\qquad c\otimes(b\otimes x)\longmapsto c\psi(b)\otimes x.
\]
因此依次沿 \(\varphi\)、\(\psi\) 扩张标量，与沿 \(\psi\circ\varphi\) 扩张标量所得的函子自然同构。
\end{proposition}
\begin{proof}
对 \(b'\in B\)，张量积的平衡关系 \(c\psi(b')\otimes b=c\otimes b'b\) 的两边都由 \(\vartheta\) 送到 \(c\psi(b'b)\)，故此公式良定义。它保持左 \(C\) 作用；对右 \(A\) 作用，有
\[
\vartheta(c\otimes b\varphi(a))
=c\psi(b)\psi(\varphi(a))
=\vartheta(c\otimes b)\cdot a.
\]
逆映射为 \(c\mapsto c\otimes1_B\)，因为
\[
c\psi(b)\otimes1_B=c\otimes b,
\qquad \vartheta(c\otimes1_B)=c.
\]
它也保持左右作用，故 \(\vartheta\) 是双模同构。

对 \(\eta_X\)，先用张量积的结合同构把源写成 \((C\otimes_B B)\otimes_A X\)，再施行 \(\vartheta\otimes_A1_X\)，便得到所列公式。其逆为
\[
c\otimes x\longmapsto c\otimes(1_B\otimes x).
\]
右端保留 \(A\) 平衡关系，因为
\[
\begin{aligned}
c\psi(\varphi(a))\otimes(1_B\otimes x)
&=c\otimes(\varphi(a)\otimes x)\\
&=c\otimes(1_B\otimes ax).
\end{aligned}
\]
上述结合同构与 \(\vartheta\) 都保留左 \(C\) 作用。对左 \(A\) 模同态 \(u:X\rightarrow X'\)，先施行 \(u\) 再施行 \(\eta_{X'}\)，或先施行 \(\eta_X\) 再施行 \(1_C\otimes_Au\)，都把 \(c\otimes(b\otimes x)\) 送到 \(c\psi(b)\otimes u(x)\)。纯张量生成源模，故得到自然性。
\end{proof}

\begin{proposition}\label{b4:prop:F-matrix-bimodule-equivalence}
设 \(k\) 为域，\(n\geq1\)，\(A=M_n(k)\)。取行向量双模 \({}_kP_A=k^{1\times n}\) 与列向量双模 \({}_AQ_k=k^{n\times1}\)，各侧作用由矩阵或标量乘法给出。行列配对与外积分别给出双模同构
\[
\begin{aligned}
\varepsilon:P\otimes_A Q&\longrightarrow k,
&p\otimes q&\longmapsto pq,\\
\theta:Q\otimes_k P&\longrightarrow A,
&q\otimes p&\longmapsto qp.
\end{aligned}
\]
第一项为 \((k,k)\) 双模同构，第二项为 \((A,A)\) 双模同构。因此 \(P:A\rightarrow k\)、\(Q:k\rightarrow A\) 给出 \(\mathbf{Bimod}\) 中对象 \(A,k\) 的等价，函子
\[
\begin{aligned}
P\otimes_A- &: A\text{-}\mathrm{Mod}\longrightarrow k\text{-}\mathrm{Mod},\\
Q\otimes_k- &: k\text{-}\mathrm{Mod}\longrightarrow A\text{-}\mathrm{Mod}
\end{aligned}
\]
互为准逆。若 \(n>1\)，则 \(A\) 与 \(k\) 不同构为环。
\end{proposition}
\begin{proof}
记 \(p_i\) 为第 \(i\) 个标准行向量，\(q_i\) 为第 \(i\) 个标准列向量，\(E_{ij}\) 为矩阵单位。矩阵乘法给出
\[
p_iE_{ab}=\delta_{ia}p_b,
\qquad E_{ab}q_j=\delta_{bj}q_a.
\]
行列配对满足 \((pa)q=p(aq)\)，因而通过 \(A\) 平衡关系，并保持两侧的 \(k\) 作用。在 \(P\otimes_A Q\) 中，
\[
\begin{aligned}
p_i\otimes q_j
&=p_iE_{ii}\otimes q_j
=p_i\otimes E_{ii}q_j
=\delta_{ij}(p_i\otimes q_i),\\
p_i\otimes q_i
&=p_i\otimes E_{i1}q_1
=p_iE_{i1}\otimes q_1
=p_1\otimes q_1.
\end{aligned}
\]
标量也可通过 \(\lambda I_n\in A\) 在张量两侧移动。展开任意行、列向量可见，\(P\otimes_A Q\) 由 \(p_1\otimes q_1\) 作为 \(k\) 向量空间生成。又有 \(\varepsilon(p_1\otimes q_1)=1\)，故 \(\varepsilon\) 为同构，其逆为 \(\lambda\mapsto\lambda p_1\otimes q_1\)。

外积满足 \((q\lambda)p=q(\lambda p)\)，所以给出 \(\theta\)。它把 \(Q\otimes_k P\) 的基 \(q_i\otimes p_j\) 送到 \(A\) 的基 \(E_{ij}\)，故为 \(k\) 线性同构。对 \(a,b\in A\)，有 \((aq)(pb)=a(qp)b\)，所以它还保持左、右 \(A\) 作用。两项同构正把复合 1态射 \(P*Q\)、\(Q*P\) 分别识别为单位双模 \(k\)、\(A\)。

对左 \(A\) 模 \(X\)，结合与单位同构给出
\[
Q\otimes_k(P\otimes_A X)
\cong(Q\otimes_k P)\otimes_A X
\cong A\otimes_A X
\cong X.
\]
对左 \(k\) 模 \(V\)，同样得到
\[
P\otimes_A(Q\otimes_k V)
\cong(P\otimes_A Q)\otimes_k V
\cong k\otimes_k V
\cong V.
\]
这两个复合的具体公式分别为
\[
q\otimes(p\otimes x)\longmapsto(qp)x,
\qquad p\otimes(q\otimes v)\longmapsto(pq)v.
\]
对 \(X\) 或 \(V\) 施行模同态，先后两种次序在纯张量上给出相同结果，因此同构自然，两个张量函子互为准逆。

当 \(n>1\) 时，\(E_{11}E_{12}=E_{12}\)，而 \(E_{12}E_{11}=0\)，故 \(A\) 非交换；域 \(k\) 交换，所以二者不能同构为环。
\end{proof}
